\centerline{\bf \Huge Тренировочная олимпиада}
\centerline{\bf \Huge }


\problem{1}
Дед Мороз катается на коньках по ледяному кольцу, двигаясь с постоянной скоростью. Одновременно из той же точки в том же направлении стартует Снегурочка. Она движется в 6,5 раз медленнее Деда. Когда Дед удаляется от Снегурочки (то есть, кратчайшее расстояние между ними растёт), она льёт слёзы, а когда приближается, то не льёт. Какое наименьшее количество целых кругов надо будет проехать Деду Морозу, чтобы к этому моменту Снегурочка залила слезами всё ледяное кольцо?

\problem{2}
Известно, что среди 63 монет 7 фальшивых. Все фальшивые монеты весят одинаково, и все настоящие монеты также весят одинаково, фальшивая монета легче настоящей. Как за три взвешивания на чашечных весах без гирь определить 7 настоящих монет?

\problem{3}
На стороне $A C$ треугольника $A B C$ отмечены точки $D$ и $E$ так, что $A D=C E$. На отрезке $B C$ выбрана точка X , а на отрезке $B D$ - точка $Y$, причём $C X=E X$ и $A Y=D Y$. Лучи $Y A$ и $X E$ пересекаются в точке $Z$. Докажите, что середина отрезка $B Z$ лежит на прямой $A E$.

\problem{4}
Найдите все вещественные положительные числа $x, y, z$, удовлетворяющие условию

$$
x+2 y+3 z=2 x y+6 y z+3 z x=3 .
$$


\problem{5}
В клетках квадрата $101 \times 101$ расставлены числа $1,0,-1$ (по одному в каждой клетке). Оказалось, что в любом клетчатом квадратике $2 \times 2$ можно выбрать три клетки так, что сумма чисел в этих клетках равна нулю. Какое наибольшее значение может принимать сумма чисел во всём квадрате?

\problem{6} На плоскости расположены $n$ квадратов $2 \times 2$ со сторонами, параллельными координатным осям. Ни один из центров этих квадратов не содержится ни в каком другом квадрате (в том числе на границе). Прямоугольник П со сторонами, параллельными координатным осям, содержит все эти квадраты. Докажите, что периметр П не меньше $4 \sqrt{n}$.